% Generated draft: requires mathematical and source review.
\documentclass[11pt]{article}
\usepackage[margin=25mm]{geometry}
\usepackage{fontspec}
\setmainfont{Latin Modern Roman}
\usepackage{amsmath,amssymb,mathtools,unicode-math}
\setmathfont{Latin Modern Math}
\usepackage{enumitem,needspace,xurl,hyperref}
\setlength{\parindent}{0pt}
\setlength{\parskip}{5pt}
\setlength{\emergencystretch}{3em}
\newcommand{\N}{\mathbb{N}}
\newcommand{\Z}{\mathbb{Z}}
\newcommand{\Q}{\mathbb{Q}}
\newcommand{\R}{\mathbb{R}}
\newcommand{\C}{\mathbb{C}}
\newcommand{\F}{\mathbb{F}}
\newcommand{\A}{\mathbb{A}}
\newcommand{\PP}{\mathbb{P}}
\newcommand{\E}{\mathbb{E}}
\newcommand{\eps}{\varepsilon}
\newcommand{\dd}{\,\mathrm{d}}
\begin{document}
\section*{UnsolvedMath ID: 30004968}
\textbf{Source problem code:} catalogue entry
\subsection{Definitions and mathematical statement}
The supplied record does not determine a well-posed mathematical problem. It introduces neither the graph parameters \(\pi\) and \(\pi^{*}\), nor the objects \(G,H,P\), their relationship, or their hypotheses. The title and supplied verification associate the fragment with nonrepetitive graph coloring, but do not specify the required definitions. The symbols \(\boxtimes\) and \(K_3\) conventionally suggest the strong graph product and the complete graph on three vertices; these interpretations alone do not recover the missing statement.

The extracted mathematical fragment is reproduced below for provenance, without asserting its validity or supplying missing assumptions:
\[
\begin{aligned}
\pi(G)
&\leq \pi(H\boxtimes P\boxtimes K_3)\\
&\leq \pi^{*}(H\boxtimes P\boxtimes K_3)\\
&\leq \pi^{*}(H\boxtimes P)\cdot 3\\
&\leq \pi^{*}(H)\cdot 4\cdot 3\\
&\leq 43\cdot 4\cdot 3=768.
\end{aligned}
\]
The final equality is inconsistent as transcribed: \(43\cdot4\cdot3=516\). Although \(4^{3}\cdot4\cdot3=768\), interpreting the supplied \(43\) as \(4^{3}\) would require confirmation and is not adopted here.

The accompanying sentence describes an unnamed problem as formerly central and open, but supplies no identifiable antecedent or question. Moreover, the supplied statement-verification note characterizes the displayed bound as proved material rather than a standalone open problem and states that no clean statement was promoted. Consequently, no quantified conjecture can be recovered from this record. Its registry label \emph{open} is supplied metadata, while its statement status is \emph{unrecoverable}; neither establishes a mathematical open-problem statement.
\subsection{Short English statement}
This record preserves an incomplete nonrepetitive graph-coloring fragment with missing definitions and hypotheses, an unidentified question, and an inconsistent numerical equality. The supplied verification describes the fragment as a proved bound. A precise open-problem statement cannot be recovered from the supplied evidence.
\subsection{Sources}
\begin{itemize}
\item Local UnsolvedMath registry, permanent ID \(30004968\), title \emph{Unspecified Nonrepetitive Graph Coloring Problem}, exact source problem code \texttt{catalogue entry}. The supplied record belongs to the Oberwolfach Reports Open Problems collection, records proposed year \(2022\), and supplies no proposer attribution. Its status is \emph{open}, its statement status is \emph{unrecoverable}, and its clean statement is absent.
\item Supplied source citation: \emph{Graph Theory (2023)}, Oberwolfach Reports, DOI 10.4171/owr/2022/1, pp.\ 5--78. Supplied source URL: \url{https://doi.org/10.4171/owr/2022/1}. The background describes Oberwolfach Reports as reports documenting mathematical workshops and their open-problem discussions. The citation's 2023 designation and the record's 2022 report/year designations require reconciliation.
\item Supplied literature source: \emph{Oberwolfach Report 1/2022}, \url{https://doi.org/10.4171/owr/2022/1}. Literature checked at: \texttt{2026-08-22}. This is supplied dated curation, not an independently verified assessment. It labels the record open while assessing the intended target as unresolved and unidentifiable because the unnamed problem and surrounding bound do not form a complete question. The background requests independent verification of time-sensitive status and cited claims.
\item Supplied statement-verification note: the official report is described as using the extracted formula to record a proved nonrepetitive-coloring bound rather than a standalone open problem; no separate clean statement is promoted, and the original record remains publicly visible for provenance. This characterization is supplied source metadata and has not been independently verified here. The report citation provides provenance, but the supplied material identifies no authoritative complete statement of the intended open problem.
\end{itemize}
\end{document}
